% Cote 119: offprints, reports, courses, CV and letters (typed and printed part): archive cote 119, site page 304 (PDF page 305).
% typescript. Draft compared: none.
% AI-transcribed by Claude Opus 5.5 (claude-opus-5-5, adaptive thinking, effort high), Message Batches API, batch r119_full (msgbatch_01DwZNm5Y7YsKBWyAN4sFs8w), 2026-11-13; not yet checked by eye.
\scanpage{119/304}
\mspage{-2-}

faisceau structural, qui trouveront leur cadre naturel quelques
années plus tard avec la théorie des topos (aboutissement natu-%
rel du point de vue faisceautique en topologie générale) [16,
SGA 4].

A partir de 1956 je suis resté en France, à l'exception
de séjours de quelques semaines ou mois dans des universités
étrangères. De 1950 à 1958 j'ai été chercheur au CNRS, avec le
grade de directeur de recherches en 1958. De 1959 à 1970 j'ai
été professeur à l'Institut des Hautes Etudes Scientifiques.
Ayant découvert à la fin de 1969 que l'IHES était subventionné
depuis trois ans par le Ministère des Armées, et après des essais
infructueux pour inciter mes collègues à une action commune sans
équivoque contre la présence de telles subventions, je quitte
l'IHES en septembre 1970.

Depuis 1959 je suis marié à une française, et je suis
père de quatre enfants. Je suis apatride depuis 1940, et ai
déposé une demande de naturalisation française au printemps 1970.

Depuis 1956 jusqu'à une date récente, mon intérêt prin-%
cipal s'est porté sur la géométrie algébrique. Mon intérêt pour
la topologie, la géométrie analytique, l'algèbre homologique ou
le langage catégorique a été constamment subordonné aux multiples
besoins d'un vaste programme de construction de la géométrie
algébrique, dont une première vision d'ensemble remonte à 1958.
Ce programme est poursuivi systématiquement dans [16, 17], d'abord
dans un isolement relatif, mais progressivement avec l'assistance
d'un nombre croissant de chercheurs de valeur. Il est loin d'être
achevé à l'heure actuelle. L'extraordinaire crise écologique que
nous aurons à affronter dans les décades qui viennent rend peu
probable qu'il le sera jamais. Elle nous imposera d'ailleurs
une perspective et des critères de valeur entièrement nouveaux,
qui réduiront à l'insignifiance (``irrelevance'') beaucoup des
plus brillants progrès scientifiques de notre siècle, dans la
mesure où ceux-ci restent étrangers au grand impératif évolution-%
niste de notre temps : celui de la survie. Cette optique s'est
imposée à moi avec une force croissante au cours de discussions
avec de nombreux collègues sur la responsabilité sociale des
scientifiques, occasionnées par ma situation à l'IHES depuis la
fin de 1969. Elle m'a conduit en juillet 1970 à m'associer à la

\margin{304} \hfill ...
